\chapter{CONCLUSION}
 
\section{Achievements}
 
This project addressed a fundamental constraint in edge deployed computer vision: the
inability of resource limited microcontrollers to process high resolution imagery without
destroying the spatial detail necessary for reliable detection. The proposed solution, a
distributed tile based inference pipeline for \gls{hpd} demonstrated that spatially
decomposed neural inference across a heterogeneous network of commodity ESP32 class
hardware is not only architecturally feasible but implementable end to end without GPU
acceleration, cloud offloading, or specialised neural processing hardware.
 
The system successfully achieved the following concrete outcomes. A complete model
training and quantisation pipeline was developed, converting a custom Improved
\gls{lbcnn} from floating-point PyTorch to fully INT8 TensorFlow Lite Micro format,
achieving approximately 4$\times$ model size reduction while preserving functional
inference capability on the target hardware. The tile extraction, attention based feature
aggregation, and binary classification pipeline was validated end to end in simulation,
with per tile inference, transmission delay, and aggregator latency each measured
independently through dedicated benchmark sketches. The distributed communication layer
was implemented using the \gls{esp}-NOW protocol, enabling sub millisecond
peer to peer transmission of 69 byte weighted feature payloads between worker nodes and
the aggregator. Fragment based tile reassembly with duplicate suppression was implemented
and validated on the ESP32-S3 worker nodes, demonstrating robustness to packet
reordering under the 250 byte \gls{esp}-NOW frame size constraint. The full quantisation
pipeline from PyTorch checkpoint through ONNX, TF SavedModel, and TFLite INT8
conversion was automated into a single reproducible script, with all three sub module
headers (tile extractor, attention, classifier) verified to contain zero float32 tensors
prior to deployment.
 
\section{Limitations}
 
Several constraints were encountered during implementation that bound the current
system's performance and generalisability.
 
\textbf{Inference latency.} The dominant bottleneck in the system is the per tile inference time on the \gls{esp}32-S3 worker nodes, measured at approximately 2.68 seconds per tile when the TFLite arena is allocated in Octal PSRAM. This far exceeds the 1,000~ms target frame period, meaning the current implementation cannot sustain real time throughput at the intended 1 \gls{fps} capture rate. The root cause is the high latency of PSRAM access relative to internal \gls{sram}: the tile extractor requires approximately 451~KB of arena memory exceeding the \gls{esp}32-S3’s available internal \gls{sram} forcing all intermediate activations through the slower PSRAM bus. This is an architectural constraint of the chosen model and hardware compatibility rather than a software defect.

 
\textbf{Model accuracy.} The reported validation accuracy of 35.87\% for the \gls{lbcnn} backbone is not good enough for deployment in a real detection system. This suggests the use of synthetic calibration data during INT8 quantisation instead of real validation imagery and a training dataset that might not be sufficiently representative of the diversity of human subjects and environmental conditions encountered in practice. The quantisation pipeline and deployment infrastructure are solid. The accuracy limitation is a function of the training data and training duration, not the architectural approach.
 
\textbf{Hardware availability.}The absence of an ESP32-CAM during development prevented the full end to end pipeline from live frame capture, tiling, distributed inference and aggregated classification to be validated on real imagery. All tile inputs to the worker nodes during benchmarking were synthetic, generated programmatically to exercise the inference and communication pathways. System behavior on real world frames with real spatial variation has not been validated.

\textbf{Fragment reassembly reliability.} The current tile fragmentation scheme sends 38
sequential \gls{esp}-NOW packets per tile with a simple ACK wait flow control. Under
network congestion or when both workers transmit simultaneously, packet collisions at the
aggregator's receive buffer can cause fragment loss and incomplete tile reassembly. While
duplicate suppression is implemented, a formal retransmission mechanism is absent.
 
\textbf{Single channel input.} The model operates on grayscale (single channel) input,
which discards chrominance information that could assist in distinguishing human subjects
from similarly shaped environmental objects in colour imagery.
 
\section{Future Work}
 
Several concrete directions emerge from this work that would substantially improve the
system's performance, robustness, and real world applicability.
 
\textbf{Reducing inference latency.} The most important step needed for the system to work in real time would be to lower the per tile inference time to under 200~ms. Three different strategies are worthy of investigation. Firstly, compression of the \gls{lbcnn} backbone using structured pruning would lower the amount of channels in each layer, thus lowering the arena size to fit into internal \gls{sram}
and removing the need to pay PSRAM latencies. Secondly, changing the \gls{lbcnn} backbone for a version of MobileNetV3(Small) or EfficientNet(Lite) optimised for TFLite Micro could provide similar feature detection capabilities while using much less memory. Lastly, the dual core structure of the ESP32-S3 could be used for running fragment reassembly and inference in a
pipeline fashion on cores~0 and~1 respectively.
\textbf{Improving model accuracy.} A validation accuracy of 35.87\% is achieved during the proof of concept training session. A systematic retraining process with an organised set of data related to human detections from datasets like INRIA Person Dataset or Crowd Human, along with appropriate augmentation, class balancing, and additional training, will be the quickest route towards achieving an accuracy score that can be implemented into a model. Moreover, using actual images instead of randomly generated Gaussian noise for INT8 quantisation will lead to better results during quantisation and less accuracy drop from Float to INT8.
\textbf{Real ESP-CAM integration.} It is most critical to validate the entire pipeline using a live ESP32 CAM stream. Doing so will provide insights into any timing issues that arise from interactions between capture, fragmentation, transmission, and inference, which cannot be identified within a synthetic benchmark test environment. Furthermore, it will indicate whether the approach of 25\%
overlap works well for tile-boundary subjects.
 
\textbf{Reliable fragment transport.} A more effective approach is using the light weighted selective retransmission mechanism whereby the worker only requests the retransmission of particular missing fragment indices instead of requesting the retransmission of all fragments for that entire tile. The other option is to use a time division multiple access technique that involves transmitting from the CAM to each worker in a certain time slot without any collisions at all.
 
\textbf{Scaling the node count.} Currently, the architecture uses two worker nodes and one aggregator. Using four or more worker nodes to handle two or three tiles each would be sufficient for processing images of greater resolution such as 448x448 using 16 tiles but without increasing the number of tiles per node. The softmax weighted average aggregation by the aggregator is independent of the number of nodes and requires only \textit{NUM\_TILES} constant.
 
\textbf{Multi class detection.} The binary human/no human classifier could be extended
to a multi class output predicting additional categories such as vehicle, animal, or
background, making the system applicable to broader surveillance and wildlife monitoring
use cases with minimal architectural change.
 
\textbf{Privacy preserving deployment.} Because all processing occurs on device with no
image data leaving the local \gls{espnow} network, the architecture is inherently
privacy preserving relative to cloud based alternatives. Future work could formalise
this property and explore the system's suitability for GDPR compliant deployments in
public spaces where raw video transmission is legally restricted.
 
\section{Final Remarks}
 
This project shows that the boundary between what is considered feasible on microcontroller class hardware and the domain of systems exclusively accelerated with \gls{gpu} is not a hard line, it is an engineering problem. In this work we show that resolution preserving HPD can be achieved without centralised compute, cloud connectivity or specialised neural processing hardware by decomposing a spatial inference task into tile level subtasks, distributing them across a network of sub five dollar microcontrollers and aggregating their outputs by means of an attention weighted fusion mechanism.
The deployed system is a research prototype with real performance limitations, most notably in inference latency and model accuracy. However, the infrastructure built the end to end quantization pipeline, the \gls{espnow} distributed communication layer, the fragment reassembly mechanism and the attention based aggregator constitutes a reusable foundation on which the identified improvements can be systematically applied. As \gls{mcu} clock speeds increase and TFLite Micro’s hardware acceleration support matures, the performance gap identified here will diminish. The architectural approach itself, spatially distributed, attention guided, feature level aggregation on heterogeneous edge nodes is not hardware specific and scales naturally to more capable devices as they become available at comparable cost and power envelopes.
The broader significance of this work lies in its demonstration that intelligent sensor systems need not be centralised to be capable. In environments where privacy, latency, connectivity, and power are all constrained simultaneously, distributed edge inference of the kind demonstrated here represents a principled and practical path forward.
